\documentclass[a4paper,12pt]{article}
\usepackage{learnpackage}


\newcommand{\mkst}{\makebox[-3pt][r]{\color{blue}\ding{78}}}
\newcommand{\lbf}[1]{\kaishu\textbf{\color{cyan}#1}}
\newcommand{\lsn}[1]{\noindent\textsc{\color{orange} #1.}}

\setlength{\headheight}{14.9998pt}
\addtolength{\topmargin}{-2.49998pt}
\pgfplotsset{compat=1.18}
%------------------Document----------------------------
\begin{document}

%-------------------------TitlePage------------------------
%\begin{titlepage}
%\end{titlepage}

%-------------------------Content------------------------
\begin{tikzpicture}[mindmap, 
	every node/.style={concept, circular drop shadow,execute at begin node=\hskip0pt}, 
	root concept/.append style={ concept color=black, fill=white, line width=1ex, text=black, font=\large\scshape}, text=white, 
	level21/.style={concept color=red,faded/.style={concept color=red!50}}, 
	level22/.style={concept color=blue,faded/.style={concept color=blue!50}}, 
	level23/.style={concept color=orange,faded/.style={concept color=orange!50}}, 
	level24/.style={concept color=green!50!black,faded/.style={concept color=green!50!black!50}}, grow cyclic, 
	level 1/.append style={level distance=4.5cm,sibling angle=90,font=\scshape}, 
	level 2/.append style={level distance=3cm,sibling angle=45,font=\scriptsize}]
	\node [root concept](Computational Complexity){波粒二象性}
	child[level21] {node {能量量子化}
	child {node {黑体与黑体辐射}}
	child [faded] {node {黑体辐射实验图像}}
	child {node {能量子}child{node{\footnotesize\(\varepsilon = h\nu\)}}}
	}
	child[level22] {node {光的粒子性}
	child {node {光电效应的实验规律}}
	child {node {爱因斯坦光电效应方程}child{node{\(E_\mathrm{km} = h\nu - W_0\)}}}
	child {node {康普顿效应}}
	child {node {德布罗意波}child{node{\footnotesize\(\lambda = \frac{h}{p}\)}}}
	}
	child [level23]{node {概率波}
	child {node {经典粒子与经典波}}
	child {node {概率波}}}
	child [level24]{node {粒子的波动性}
	child {node {光的波粒二象性}}
	child {node {光的波动性与粒子性的统一}}
	child {node {物质波}}};

	\begin{pgfonlayer}{background} 
		\colorlet{upperleft}{green!50!black!25} 
		\colorlet{upperright}{orange!25} 
		\colorlet{lowerleft}{red!25} 
		\colorlet{lowerright}{blue!25}
		\fill [upperleft] (Computational Complexity) rectangle ++(-.5\textwidth,10); 
		\fill [upperright] (Computational Complexity) rectangle ++(.5\textwidth,10); 
		\fill [lowerleft] (Computational Complexity) rectangle ++(-.5\textwidth,-10); 
		\fill [lowerright] (Computational Complexity) rectangle ++(.5\textwidth,-10);

		\shade [left color=upperleft,right color=upperright] ([xshift=-1cm]Computational Complexity) rectangle ++(2,10); 
		\shade [left color=lowerleft,right color=lowerright] ([xshift=-1cm]Computational Complexity) rectangle ++(2,-10); 
		\shade [top color=upperleft,bottom color=lowerleft] ([yshift=-1cm]Computational Complexity) rectangle ++(-10,2); 
		\shade [top color=upperright,bottom color=lowerright] ([yshift=-1cm]Computational Complexity) rectangle ++(10,2); 
	\end{pgfonlayer} 
	\end{tikzpicture}
\end{document}
